\scanpage{431}
\mspage{218}

it is likewise \unsure{for} the action induced on each $\pi(I')$ $(I'\subset I=I(\pi))$ [it suffices to take the $\pi(I')$ for $I'$ of the form $I\setminus\{j\}$, $j\in I$] --- \uuline{and} that finally the action of $\Gamma^{\natural}$ \deleted{relative to \unsure{an} $i\in I$} on $\pi'(I)$ (deduced from $\pi'$ by plugging the holes at $I$, so that $\pi'(I)$ has a single class of loops; \unsure{$\pi'(I)$} is also deduced from the \inserted{actual} action of $\Gamma$ on $\pi(I)$ --- with $0$ \inserted{classes of} loops --- by making ``a hole'' corresponding) relative to \inserted{a} point $i\in I$, normalizes no loop subgroups\ldots\ But I presume that the first condition (action \inserted{of $\Gamma^{\natural}$} normalizing none of the $L_i$) implies the others.

\margin{at least if $g\ne 0$}

But one must \inserted{in the definition of admissibility} also take into account condition $(4^\circ)$, \deleted{\illegible{1}}\inserted{which implies the }existence of enough liftings\ldots\ .

I come (painfully!) to a

\uuline{Provisional profinite conjecture} For $\Gamma$ a given profinite group, an outer action on an arithmetized anabelian $(\pi,a)$ of type $(g,\nu)$ is called \uline{admissible}, if

\margin{Hoaxed cf. below\ldots}

a) \deleted{\illegible{1}} The homomorphism $\Gamma\longrightarrow\boldsymbol{\Gamma}_a$ has open image

b) The \uline{actual} \unsure{actions} deduced by pluggings \unsure{of|of the} holes $(i\in I'\subset I=I(\pi))$ normalize no loop subgroups.

\scanpage{432}

c) There exist infinitely many \inserted{$\pi$-conjugacy classes of the} germs \deleted{\illegible{1}} splittings of the extension of $\Gamma$ by $\pi$, which normalize no loop subgroup of $\pi$.

\margin{This \inserted{condition c)} probably excludes cases like $\Gamma=\unsure{\mathcal N_{g,\nu}|N_{g,\nu}}$ with $(g,\nu)\ne(0,3)$, for the extension $\unsure{\mathcal M_{g,\nu}|M_{g,\nu}}$ of $\unsure{\mathcal N_{g,\nu}|N_{g,\nu}}$ by $\unsure{\widehat{\pi}_{g,\nu}|\widehat{\Pi}_{g,\nu}}$ probably admits no germs of \illegible{1} ($\overset{?}{\Longleftrightarrow}$ $\unsure{\mathcal U_{g,\nu}|U_{g,\nu}}$ over $\unsure{\mathcal M_{g,\nu}|\mathcal N_{g,\nu}}$ \unsure{has|admits} probably no \unsure{étale multisection|étale multisections}\ldots)}

An actual action of $\Gamma$ on $\pi$ is called actual, \deleted{and} if the outer action is actual, and if it normalizes no loop subgroup.

This being so:

a) If $\Gamma$ operates actually in an admissible way, we have
\[
\pi^\Gamma=(1)
\]
(perhaps even $\pi^{\unsure{\Gamma^+|\Gamma^\dagger}}=(1)$) \deleted{I also presume that the only \unsure{sub-}}

\deleted{b) If $\Gamma$ operates by outer automorphisms in an admissible way, \deleted{then} leaving $I'\subset I(\pi)$ stable, then ditto on $\pi(I')$ --- likewise for admissible actual actions.}

\margin{\deleted{\illegible{1}}}

\deleted{c) Under the preceding conditions,}

\scanpage{433}
\mspage{219}

\deleted{the canonical map of $\Gamma$-sets
\[
\mathrm{Pt}(B_{\pi,\Gamma^{\natural}})\longrightarrow
\mathrm{Pt}(B_{\pi(I'),\Gamma^{\natural}})
\]}

b) If $i\in I$ is stable under $\Gamma$ operating \deleted{actually}\inserted{by outer automorphisms} in an admissible way, then the actual action on $\pi'=\pi(i)$ normalizes no $L'_j$ $(j\in I\setminus\{i\})$.

c) If $\Gamma$ operates actually on $\pi$ in an admissible way, \inserted{\unsure{while} leaving $i\in I$ fixed,} then the corresponding actual action (\uline{by} \deleted{\illegible{3}}\inserted{passage to the quotient}) on $\pi^{\prime\natural}=\unsure{\pi(i)|R(i)}$ normalizes no $L'_j$ $(j\in I\setminus\{i\})$.

d) If $\Gamma$ \deleted{operates actually on $\pi$ in a way}\inserted{operates actually on $\pi$ in an} admissible way, then the map
\[
\unsure{\beta}\qquad
\mathrm{Pt}_{\mathrm{ad}}(B_{\pi,\Gamma^{\natural}})\longrightarrow
\mathrm{Pt}_{\mathrm{ad}}(B_{\pi',\Gamma^{\natural}})
\]
defined by c) is \uline{injective}, and \inserted{the complement of} its image is \deleted{finite} equal to $\{P\}$, where $P$ is defined by the ``hole'' $i$.

But this last part of the assertion, going beyond mere \uline{injectivity} --- and characterizing the image, now seems to me quite doubtful ---

\scanpage{434}

indeed, it suffices to look at a parameter scheme $S$, of finite type over \unsure{$Q$|$\mathbb{Q}$} (a $K(\pi,1)$ preferably, for example an \unsure{algebraic} curve) and to take for $\Gamma$ (not the fundamental group of its generic point, \inserted{i.e. of a field,} \unsure{but)} the fundamental group of $S$ itself. The datum of an admissible outer action of $\Gamma$ on a $\pi$ \unsure{corresponds} morally to that of a family of curves $U$ of type $g,\nu$ parameterized by $S$. Condition c) is satisfied for example if $U$ comes from an algebraic curve over the base field itself, so no problem -- and it is plain that surjectivity goes haywire. The difficulty comes from the fact that two \uline{distinct} sections of $U$ over $S$ may not be \uline{disjoint} !

\margin{The splittings of extensions correspond \unsure{to the} sections of $U$ over $S$}

\scanpage{435}
\mspage{220}

One must be able to speak of sections ``\unsure{strictly|strictly} distinct'' i.e. distinct at every point of $S$ -- which, in profinite translation, \unsure{amounts} to comparing two splittings of the extension of $\Gamma$ by $\pi$, with \deleted{\unsure{two}}\inserted{the} \inserted{germs of} splittings over \unsure{$\boldsymbol{\Gamma}'_{\cdot}\subset\boldsymbol{\Gamma}$|$\boldsymbol{\Gamma}'\subset\boldsymbol{\Gamma}$} of $\Gamma$ (as \unsure{an extension} of $\boldsymbol{\Gamma}_Q$ by a geometric part) which correspond to the \unsure{points} of $S$ -- i.e. precisely the \uline{admissible} (germs of splittings). One has the impression of going around in a vicious circle or almost -- it would seem that one should not try to swallow too much at once -- deal from the outset with \uline{all} profinite groups $\Gamma$ at once -- but rather restrict \unsure{oneself} \inserted{first} to those which, morally, correspond to \inserted{\unsure{$\pi_1$|$\Pi_1$} of} schemes of finite type over \unsure{$Q$|$\mathbb{Q}$} -- some $K(\pi,1)$ let us say, or \inserted{even} varieties

\scanpage{436}

elementary in Artin's sense -- and in the definitions resp. conjectures, pull oneself up by one's shoelaces, by inducting on the dimension.

At the first notch then (dimension $0$) one would restrict oneself to actions of groups $\Gamma$ which are (modulo at least passage to an open subgroup) ``arithmetically faithful'' \unsure{for example|for example} $\Gamma\to\boldsymbol{\Gamma}_a$ injective.

\unsure{In|\uline{In}} that case, the conjecture as stated just now seems reasonable, or at least not necessarily crazy. The decisive test\unsure{,|,:} it is true, would be the possibility of obtaining the ``curves'' thus defined as coverings of $\mathbb{P}^1\setminus\{0,1,\infty\}$ \ldots

\scanpage{437}
\mspage{221}

\uline{33} Topological digression :
Anti-involutions of oriented \unsure{algebroid} surfaces

\margin{In the end I look only at compact surfaces -- Yet the noncompact case would also be very interesting to look at -- to try to recover by an ``algebraic'' route, from the outer automorphism of order $2$ in $\unsure{\pi_1}$, the arrangement of the curve's ``points at infinity'' on the connected components of \deleted{\illegible{1}} $X^\sigma$ $(X=\widehat U)$}

\unsure{We call} an \unsure{algebroid} surface a surface $U$ of the form $\widehat X\setminus S$, $X$ a compact orientable \inserted{surface}, $S$ finite. Then $X$ is determined up to unique homeomorphism as the ``pure compactification'' of $U$, \unsure{denote it|we denote it} $\widehat U$. \deleted{\unsure{The}} The homeomorphisms of $U$ onto itself are identified with the morphisms of $X$ which map $S$ into itself. The orientations of $X$ are in $1$-$1$ correspondence with those of $U$. The anti-involutions of $U$ (\unsure{assumed} oriented) are in $1$-$1$ correspondence with those of $X$ which preserve $S$.

We then find an equivalence between the category of oriented \unsure{algebroid} surfaces \inserted{$U$,} equipped with an anti-involution $\sigma$, and of compact surfaces with boundary $Y$, equipped with a finite subset $T$, to $(U,\sigma)$ corresponds $(\widehat U/\sigma,(\widehat U\setminus U)/\sigma)$ -- and to $(Y,T)$ corresponds $\widetilde Y-\widetilde T$, where $\widetilde Y$ is the ``oriented double'' of $Y$ \inserted{(which is a \unsure{compact} oriented surface)}, and $\widetilde T$ the inverse image of $T$ in $\widetilde Y$.

We are now interested in the

\scanpage{438}

case where $U$ is connected of type $(g,\nu)$, first examining the case $\nu=0$, \inserted{i.e. $U=\widehat U$, $S=\varnothing$} (to which the general case \unsure{will reduce}). \deleted{This} \inserted{The condition $\nu=0$} corresponds to the case where $T=\varnothing$ -- thus the \inserted{$U=X$} \deleted{$X$} considered \deleted{are} are identified with the oriented doublings of certain compact manifolds with boundary $Y$. Which ones ? Obviously $Y$ must be \uline{connected}, and \uuline{nonorientable}.

\margin{it is not true that $Y$ is necessarily nonorientable -- only in the case where $X^\sigma=\varnothing$}

So $Y$ is \unsure{characterized}, up to homeomorphism, by its genus $\gamma$ \inserted{and the number $\mu$ of connected components \inserted{of the boundary}} -- it can be obtained from the projective plane \inserted{$Y_0$ over the reals} by cutting out $\gamma$ disjoint disks in it, and gluing Möbius strips back in, \inserted{then \unsure{removing again} $\mu$ open disks,} -- which gives
\[
\begin{aligned}
\chi(Y)&=\underset{\substack{\text{''}\\1}}{\chi(Y_0)}
-\deleted{\unsure{\chi}}\unsure{\gamma}\cdot
\underbrace{\chi_!(\text{open disk})}_{+1}
+\unsure{\gamma}\chi_!(\text{open Möbius strip})\\
&\hspace{11em}\unsure{+|-}\mu\underbrace{\chi_!(\text{open disk})}_{+1}
\end{aligned}
\]
now the \inserted{open} Möbius strip (obtained from $Y_0$ by \unsure{removing} a \inserted{closed} disk) has \unsure{a|a} $\chi_!$ equal to :
\[
\underset{\substack{\text{''}\\1}}{\chi(Y_0)}-
\underset{\substack{\text{''}\\1}}{\chi(\text{closed disk})},
\quad \text{i.e. }0,\quad \text{hence}
\]
\[
(1)\qquad \chi(Y)=1-(\unsure{\gamma}+\mu)
\]
\margin{$Y$ compact connected nonorientable surface without boundary of genus $\gamma$}

On the other hand, we \unsure{will have|have}
\[
\chi(X)=\chi(X\deleted{\unsure{)}}\setminus X^\sigma)
\]
for $X^\sigma$ is a \unsure{union} of circles so its $\chi$ is zero,

\scanpage{439}
\mspage{222}

Now, since $X\setminus X^\sigma$ is an étale covering of order $2$ of $\unsure{\operatorname{Int}Y}$, we have
\[
\chi(X\setminus X^\sigma)=2\chi(\operatorname{Int}(Y))
\]
finally
\[
\chi(Y)=\chi(\underline{Y\setminus\partial Y})=\chi(\operatorname{Int}Y)
\]
for the same reason as \unsure{before} $(\chi(\partial Y)=0)$,
hence finally
\[
(2)\qquad \chi(X)=2\chi(Y)=2(1-(\gamma+\mu))
\]
\deleted{\illegible{1}} or again, since $\chi(X)=2-2g$
\[
(3)\qquad g=\gamma+\mu
\]

Thus we find

\margin{\illegible{1}}

\uline{Proposition} The category of oriented \inserted{connected compact} surfaces of genus $g$ equipped with an anti-involution $\sigma$ is equivalent to that of \uline{nonorientable} compact manifolds with boundary, of type $(\gamma,\mu)$ ($\gamma$ the genus, $\mu$ the number of holes), with $\gamma+\mu=g$.

\margin{no, we find only a full subcategory of that of all the $(X,\sigma)$. On the other hand there are also the oriented doublings $Y\amalg_{\partial Y}Y$ of compact \uline{orientable} surfaces with \uline{nonempty boundary} -- if $Y$ is of type $(\gamma,\mu)$, $X$ is of genus $g$ \unsure{with} $g=2\gamma+\mu-1$ (where $\gamma\geq0$, $\mu\geq1$)}

\uline{Corollary} There are exactly $g+1$ isomorphism classes of systems $(X,\sigma)$, classified by $0\leq\mu\leq g$, where $\mu=\operatorname{card}\unsure{\pi_0|\Pi_0}(X^\sigma)$.

\margin{This only works by restricting oneself to the $(X,\sigma)$ such that $(X/\sigma)$ nonorientable -- \unsure{cf.|cf.} opposite \deleted{\unsure{for}} the \deleted{\unsure{general}} case \inserted{$X/\sigma$ orientable}}

$X$ \inserted{$=X_g$} is a compact oriented surface of

\deleted{Anti-involutions of \unsure{oriented} surfaces of \deleted{\illegible{1}} \unsure{algebroid} type (i.e. of the form $X\setminus S$, with $X$ compact, $S$ finite). We \deleted{will admit} \inserted{know} that they are classified, modulo isomorphism \inserted{up to}, by the systems}

\scanpage{440}

genus $g$, this also means that in the group $A_g=\mathrm{Aut}(X_g)$, there are exactly $g+1$ conjugacy classes of elements $\sigma$ satisfying
\[
(4)\qquad \sigma^2=1,\qquad \mathrm{sg}(\sigma)=+1.
\]
(i.e. which are anti-involutions). They give $g+1$ conjugacy classes of elements of $A_g/A_g^\circ=\unsure{\mathcal{T}_g}$, satisfying the \unsure{same|same} relations. We will show (we hope!) that these classes are distinct, and that every $\sigma\in\unsure{\mathcal{T}_g}$ satisfying relations (4) is in one of these classes.

We will admit the

\uline{Lemma} Every anti-involution of the \inserted{unit} disk \inserted{(\unsure{case|or} of $\mathbb C$)} is conjugate to $z\mapsto\overline z$, and consequently \unsure{the set} of its fixed points is homeomorphic to $\mathbb R$, and in particular is \uline{connected}.

\uline{Theorem} Let $U$ be an oriented surface \inserted{\unsure{(Hausdorff connected paracompact)}}, \inserted{\uline{not isomorphic to} $S^2$,} $\sigma$ an anti-involution of $U$, $\widetilde U$ a universal covering of $U$, whence $\pi=\mathrm{Aut}(\widetilde U/U)$ and an extension \unsure{$E$|$\mathcal E$} of \unsure{$\mathbb Z/2$} by $\pi$, formed
